\section{Exact resource counts and their optimization}\label{res:section}

The number of cone factors, the cone dimension, the ambient and
restricted barrier parameters, and the number of free affine coordinates
measure different aspects of a formulation. This section determines exact
integer tradeoffs among them. Each optimization statement concerns only
the resource model it names and does not by itself bound running time.
The lifts and slack factorizations are those of
Sections~\ref{sec:products}--\ref{sec:concrete-barriers}, interpreted in
the factorization framework of \citet{GPT2013}; the balancing and integer
optimization arguments are elementary consequences of the curvature and
nullity bounds proved here.

\subsection{Real PSD order caps}
Let $N\geq4$ and $R\geq2$ be integers, and consider exact lifts of
$\B_2^N$ through real PSD blocks of orders $2\leq r_i\leq R$ with globally
bi-$C^1$ full-slack contact selections. Put
\[
 c_r=\lfloor r^2/4\rfloor,\qquad d_r=r(r+1)/2,\qquad
 M=\sum_i d_{r_i},\qquad \rho=\sum_i r_i.
\]
Here $c_r$ is the capacity of an order-$r$ block. Optional ray factors
add resources but no curvature, so we omit them from the lower bounds.
The ambient product cone has optimal barrier parameter $\rho$, even among
coupled barriers; this does not determine the parameter after the affine
equations are imposed.
By Corollary~\ref{cor:ball-global-saturation},
\begin{equation}\label{res:capacity}
 \sum_i c_{r_i}\geq N.
\end{equation}
Indeed, the local lower bound is $N-1$, and equality would require the
single matching spin factor $Q_{N+1}$, which is not among these real PSD
cones when $N\geq4$.

\begin{proposition}[Exact integer consequences]\label{res:integer}
Define $K_R(N),V_R(N),F_R(N)$ as the minima of, respectively,
$\sum_r a_r$, $\sum_r ra_r$, and $\sum_r d_ra_r$ over nonnegative
integers $a_2,\ldots,a_R$ satisfying $\sum_r c_ra_r\geq N$.
Every lift in the stated class, with $L$ PSD blocks, satisfies
\[
 L\geq K_R(N),\qquad \rho\geq V_R(N),\qquad M\geq F_R(N).
\]
Write $N=qc_R+t$ with $0\leq t<c_R$. Then
\begin{align}
 K_R(N)&=\lceil N/c_R\rceil,\label{res:K}\\
 V_R(N)&=qR+\begin{cases}0,&t=0,\\\lceil2\sqrt t\rceil,&t>0.\end{cases}
 \label{res:V}\\
 F_R(0)&=0,\qquad
 F_R(n)=\min_{2\leq r\leq R}\{d_r+F_R((n-c_r)_+)\}.
 \label{res:F}
\end{align}
In particular,
\begin{align}
 \rho&\geq\max\{\lceil RN/c_R\rceil,\lceil2\sqrt N\rceil\},\nonumber\\
 M&\geq\max\{\lceil(2+1/\lfloor R/2\rfloor)N\rceil,
                       2N+\lceil\sqrt N\rceil\}.\label{res:ratios}
\end{align}
Two explicit cases are
\[
 F_4(N)=10\lfloor N/4\rfloor+3(N\bmod4),\qquad
 F_5(N)=\lceil5N/2\rceil\quad(N\geq4).
\]
These values are exact for the integer problems; for $R\geq4$, we do not
claim that lifts attain them.
\end{proposition}
\begin{proof}
Formula~\eqref{res:K} is immediate. Extend $c_r$ to nonnegative integer $r$.
Its increments $c_r-c_{r-1}=\lfloor r/2\rfloor$ are nondecreasing.
Moving one unit from a smaller positive order to a larger order below $R$
therefore increases or preserves capacity. Repeating, a list of total
order $S=\ell R+j$, $0\leq j<R$, has capacity at most
$\ell c_R+c_j$. This relaxation permits zero and one as orders, which
only strengthens its use as an upper bound on capacity. The smallest
order with $c_r\geq t>0$ is $\lceil2\sqrt t\rceil$. Maximal-order blocks
and this last block attain~\eqref{res:V}; the capacity bound just proved
excludes every smaller total order. Selecting the first block of an
optimal multiset proves both inequalities in~\eqref{res:F}.

For $r=2j$ or $2j+1$ one has $d_r/c_r=2+1/j$, whereas $c_r/r$
increases with $r$. These give the cap-dependent bounds in
\eqref{res:ratios}. Also
$N\leq\sum_i r_i^2/4\leq\rho^2/4$ and
$d_r=2c_r+\lceil r/2\rceil$, proving the other two bounds.
For $R=4$, orders two and three cost three per unit of capacity;
order four costs ten for four units. Using the maximal possible number
of complete four-unit blocks and paying three per remaining unit is
optimal; overshooting by another four-unit block costs ten rather than
at most nine. For $R=5$, the ratio lower bound is $5/2$. Every even
integer at least four is a nonnegative combination of four and six,
realized by orders four and five at that ratio. For odd $N\geq5$, add
one order-two block to a realization of $N-1$. This attains the ceiling.
\end{proof}

\subsection{Grouped Schur constructions and exact minima for small orders}
Partition $[N]$ into nonempty groups $G_j$ of sizes $g_j\leq R-1$.
The formulation
\begin{equation}\label{res:schur}
 \begin{pmatrix}s_j&x_{G_j}^T\\x_{G_j}&I_{g_j}\end{pmatrix}\succeq0,
 \qquad \sum_{j=1}^k s_j=1
\end{equation}
projects exactly onto $\B_2^N$. At a boundary point the allocation is
$s_j=\norm{x_{G_j}}^2$. Polynomial slack factors are
\[
 A_j(x)=\begin{pmatrix}\norm{x_{G_j}}^2&x_{G_j}^T\\x_{G_j}&I\end{pmatrix},
 \qquad B_j(y)=\tfrac12(1,-y_{G_j})^T(1,-y_{G_j}),
\]
since $\sum_j\tr(A_j(x)B_j(y))=\tfrac12\norm{x-y}^2=1-x^Ty$
on the unit sphere. Thus this construction meets the global hypotheses.
For a prescribed $\lceil N/(R-1)\rceil\leq k\leq N$, write
$N=ak+t$ with $0\leq t<k$. Balancing the group sizes minimizes the
ambient dimension within this construction and gives
\begin{equation}\label{res:schur-ledger}
 (L,\rho,M,\nu_{\rm slice})
 =\bigl(k,N+k,(k-t)d_{a+1}+td_{a+2},k\bigr).
\end{equation}
Here $\nu_{\rm slice}=k$ is also the intrinsic barrier parameter of this
fixed affine domain, by Theorem~\ref{thm:grouped-intrinsic}. Strict convexity
of $d_{g+1}$ in $g$ proves balancing. The minimum count is
$\lceil N/(R-1)\rceil$, whereas the minimum dimension is $3N$:
\[
 d_{g+1}-3g=(g-1)(g-2)/2\geq0,
\]
with equality for groups of size one or two. Eliminating the constant
identity blocks leaves $N+k$ variables and one equality. In the unreduced cone encoding
there are $1+\sum_jd_{g_j}$ affine equations, giving the same affine
dimension $N+k-1$.

The lower bounds above and this construction give the following exact
minima, attained simultaneously by one lift, in the global
contact-selection class:
\begin{equation}\label{res:small-real}
 \begin{array}{c|ccc}
 \text{real PSD dictionary}&L_{\min}&M_{\min}&\rho_{\min}\\\hline
 r\leq2&N&3N&2N\\
 r\leq3&\lceil N/2\rceil&3N&\lceil3N/2\rceil.
 \end{array}
\end{equation}
For the fixed complex order-two dictionary, the corresponding triple is
\begin{equation}\label{res:small-hermitian}
 (\lceil N/2\rceil,4\lceil N/2\rceil,2\lceil N/2\rceil)
 \quad(N\geq4);
\end{equation}
for the fixed quaternionic order-two dictionary it is
\[
 (\lceil N/4\rceil,6\lceil N/4\rceil,2\lceil N/4\rceil)
 \quad(N\geq6).
\]
The exceptions are $(1,4,2)$ for
complex $N=3$ and $(1,6,2)$ for quaternionic $N=5$. Indeed, these cones
are $Q_4,Q_6$, with capacities two and four; the claims follow from the
strict capacity increment and the fixed-tail construction of
Section~\ref{sec:orbits}. Dimensions count full dictionary cones, even
for a partially filled final group.

\subsection{Tradeoffs for packing columns into shared PSD blocks}
Let $C=(\B_2^s)^b$, where $b\geq2$ and $s\geq3$ are integers. Choose an integer
$1\leq p\leq s$ and split each source into $h=\lceil s/p\rceil$
nonempty groups of at most $p$ coordinates. Pad columns by zeros and
pack the $H=bh$ columns into $\ell$ nonempty blocks. Block $j$ has a
$p\times c_j$ matrix $W_j$ and a free symmetric matrix $S_j$:
\begin{equation}\label{res:packed-domain}
 \begin{pmatrix}S_j&W_j^T\\W_j&I_p\end{pmatrix}\succeq0,
 \qquad \sum_{\gamma\text{ belonging to }a}(S_{j(\gamma)})_{\gamma\gamma}=1
 \quad(a=1,\ldots,b).
\end{equation}
Every column belongs to one source group and appears once.

\begin{proposition}[Exact resource counts of finite packings]\label{res:packing}
Write $H=q\ell+t$ with $0\leq t<\ell$. For fixed $p,\ell$, balanced column
counts minimize the maximum order, the ambient cone dimension $M$, and the
number $V$ of free coordinates, giving
\begin{align}
 L&=\ell,&r_{\max}&=p+\lceil H/\ell\rceil,&\rho&=p\ell+H,\nonumber\\
 M&=(\ell-t)d_{p+q}+td_{p+q+1},&
 V&=bs+(\ell-t)d_q+td_{q+1},&\nu_{\rm slice}&=H.
 \label{res:packing-ledger}
\end{align}
The last value is the exact parameter of the restricted standard barrier; in fact, the fixed affine domain has intrinsic parameter
$H$. The count $V$ is taken before the $b$ independent allocation
equations are imposed. Under an order cap $R$, all undominated resource
vectors of this fixed-width family are obtained by enumerating
\begin{equation}\label{res:enumerate}
 1\leq p\leq\min(s,R-1),\qquad
 \lceil H/(R-p)\rceil\leq\ell\leq H
\end{equation}
and discarding dominated vectors.
\end{proposition}
\begin{proof}
Schur complementation gives $D_j=S_j-W_j^TW_j\succeq0$.
The allocation identities bound the squared norm of each source vector
by one. Conversely, any source points in the balls admit diagonal residuals with
the required nonnegative sums. This proves exactness.
Moving a column from a block with at least two more columns to a smaller
block decreases both sums of $d_{p+c}$ and of $d_c$ and does not
increase the maximum order. This proves balancing and all counts except the barrier equality.

The barrier $-\sum_j\log\det D_j$ has parameter at most $\sum_jc_j=H$
by the matrix-epigraph calculation in Theorem~\ref{thm:packing-intrinsic}.
Project the closed slice onto all $W_j$ and all diagonal entries of $S_j$.
Its image is exactly the grouped domain of
Theorem~\ref{thm:grouped-intrinsic}, with $H$ nonempty groups. Fibers are
bounded: $|D_{uv}|^2\leq D_{uu}D_{vv}$ and each diagonal residual is at
most one. There is a strictly feasible point at $W=0$ and positive
allocation to every column. Lemma~\ref{lem:bounded-fiber-barrier}
therefore gives an intrinsic lower bound $H$, proving equality.
Finally, balanced maximum order is at most $R$ exactly in
\eqref{res:enumerate}; widths $p>s$ only add zero padding and resources.
\end{proof}

For example, the exact minimum factor count and ambient rank in this
family are
\[
 \min_p\left\lceil\frac{H_p}{R-p}\right\rceil,
 \qquad
 \min_p\left\{H_p+p\left\lceil\frac{H_p}{R-p}\right\rceil\right\},
 \qquad H_p=b\lceil s/p\rceil.
\]
Minimizing $M$, or a combined cost such as $M\sqrt{H_p}$, requires the
finite enumeration over $\ell$; a continuous balanced aspect ratio would
ignore the overshoot of the final block. The answer changes
if only free affine coordinates are counted. For every $\theta\geq0$, put
$p_*=\min(s,R-1)$ and $H_*=b\lceil s/p_*\rceil$. Then
\begin{equation}\label{res:free-optimum}
 \min V\nu_{\rm slice}^{\theta}=(bs+H_*)H_*^{\theta}.
\end{equation}
Indeed splitting a column block of sizes $c_1+c_2$ saves $c_1c_2$ free
coordinates, so for fixed $p$ the unique optimal column allocation is
$\ell=H_p$. The function $(bs+H)H^\theta$ increases with $H>0$.
Thus the largest feasible $p$ is optimal; a smaller width with the same
value of $\lceil s/p\rceil$ ties it. With one column per block, the
optimizers are the unshared grouped Schur blocks.

Without a cap, direct Lorentz lifting has
$(L,\rho,M,\nu_{\rm slice})=(b,2b,b(s+1),b)$.
Every packing instead satisfies
\begin{equation}\label{res:packing-dimension}
 M\geq(2p+1)H_p\geq b(2s+1),\qquad H_p\geq b,
\end{equation}
since
$d_{p+c}=(2p+1)c+(p-c)(p-c+1)/2$ and the last term is nonnegative
for integer $p-c$. Thus the direct Lorentz formulation has strictly
smaller cone dimension and no larger restricted parameter, with equality
for $p=s$. Nevertheless, a packing can use
a single PSD factor, of minimum order within this family
\[
 r_{\rm one}=\min_{1\leq p\leq s}\{p+b\lceil s/p\rceil\},
 \qquad r_{\rm one}<2b\ \Longleftrightarrow\ b>s.
\]
If $\lceil s/p\rceil\geq2$, the order exceeds $2b$; otherwise $p=s$
and the order is $s+b$. This proves the equivalence.

Now cap both the Lorentz dimension and the PSD order by an integer
$3\leq R<s+1$. Grouped Lorentz blocks use $h_L=\lceil s/(R-2)\rceil$ groups per source and have
\[
 (L,\rho,M,\nu_{\rm slice},V)
 =(bh_L,2bh_L,bs+2bh_L,bh_L,bs+bh_L).
\]
Grouped PSD Schur blocks instead use $h_S=\lceil s/(R-1)\rceil$,
with $(L,\rho,\nu_{\rm slice},V)=(bh_S,bs+bh_S,bh_S,bs+bh_S)$
and balanced dimension from~\eqref{res:schur-ledger}, summed over sources.
The PSD blocks have no larger free-coordinate and restricted-parameter
counts, although a PSD cone of order $R$ has $d_R$ scalar coordinates and
a Lorentz cone of dimension $R$ only $R$. Under a cap of $D$ scalar
coordinates per cone, the PSD order is only about $\sqrt{2D}$. When
source-boundary and last-block effects are negligible, grouped Lorentz has
$(L,\nu_{\rm slice},\rho,M)\sim(bs/D,bs/D,2bs/D,bs)$,
whereas balanced PSD packing has
$(2bs/D,\sqrt2bs/\sqrt D,2\sqrt2bs/\sqrt D,2bs)$.
In this comparison, the advantage of packing comes from sharing blocks
across sources, not from using fewer scalar coordinates.

\subsection{Maximal capacity for a nullity budget and a general resource inequality}
Return to globally bi-$C^1$ full product-slack factorizations over $L$
real PSD blocks of orders at most $R$. Write $n=b(s-1)$ and let
$\nu=\nustd$. At a simultaneous contact, let $c_i=\operatorname{nullity}X_i$.
The mixed derivative has rank $n$, so
\begin{equation}\label{res:nullity-capacity}
 n\leq\sum_i c_i(r_i-c_i),\qquad \sum_i c_i\leq\nu.
\end{equation}
The second inequality is Lemma~\ref{lem:boundary-order}.
For integers $z\geq0$, define the largest capacity allowed by a nullity
budget $z$ as
\[
 {\cal C}_R(L,z)=\max\left\{\sum_i c_i(R-c_i):
         0\leq c_i\leq R,\ c_i\in\mathbb Z,\ \sum_i c_i\leq z\right\}.
\]
Put $m=\lfloor R/2\rfloor$ and $z'=\min(z,Lm)$, and write $z'=qL+t$ with
$0\leq t<L$. Concavity and balancing give the exact expression
\begin{equation}\label{res:nullity-envelope}
 {\cal C}_R(L,z)=(L-t)q(R-q)+t(q+1)(R-q-1).
\end{equation}
Indeed, replacing $c_i>m$ by $R-c_i$ preserves capacity and uses less
budget, and below $m$ all marginal increments are positive; so using the
budget up to $Lm$ and balancing is optimal. Together with the sourcewise
bound on selected ranks in Theorem~\ref{thm:smooth-product-rank}, this
proves the joint necessary conditions
\begin{equation}\label{res:joint-conditions}
 \lfloor\nu\rfloor\geq b\lceil s/(R-1)\rceil,\qquad
 {\cal C}_R(L,\lfloor\nu\rfloor)\geq n.
\end{equation}
Grouped Schur lifts attain the first bound. We do not claim that every
integer point satisfying these conditions is realized by a lift. If the
total nullity is $z\leq LR/2$, Cauchy--Schwarz also gives
$n\leq Rz-z^2/L$. Here $n>0$ forces $z>0$ and
$Rz-n\geq z^2/L>0$, so $L\geq z^2/(Rz-n)$.

\begin{theorem}[Resource product bound]\label{res:holder}
Let $J_\omega=\sum_i r_i^\omega$. Under the preceding contact and
standard-barrier hypotheses, for $\theta\geq0$ and
$0\leq\omega\leq2+\theta$,
\begin{equation}\label{res:holder-bound}
 J_\omega\nu^\theta\geq
 C_\theta\frac{n^{1+\theta}}{R^{2+\theta-\omega}},
 \qquad C_\theta=\frac{(2+\theta)^{2+\theta}}
                         {(1+\theta)^{1+\theta}}.
\end{equation}
In the continuous relaxation of the capacity bound, equality requires
$p:c=(1+\theta):1$ in each block, where $p=r-c$.
\end{theorem}
\begin{proof}
For $c_i>0$, write $p_i=r_i-c_i=x_i r_i$. Maximizing
$x(1-x)^{1/(1+\theta)}$ on $[0,1]$ gives
\[
 p_ic_i\leq K_\theta R^{(2+\theta-\omega)/(1+\theta)}
 r_i^{\omega/(1+\theta)}c_i^{\theta/(1+\theta)},\qquad
 K_\theta=\frac{1+\theta}{2+\theta}
             (2+\theta)^{-1/(1+\theta)}.
\]
Zero-capacity terms can be discarded. For $\theta>0$, sum and apply
H\"older with exponents $1+\theta$ and $(1+\theta)/\theta$.
Use~\eqref{res:nullity-capacity} and $\sum c_i\leq\nu$, then raise to
power $1+\theta$. Since $K_\theta^{-(1+\theta)}=C_\theta$, the result
follows. At $\theta=0$ it is simply the sum of
$p_ic_i\leq R^{2-\omega}r_i^\omega/4$.
The shape maximum is unique at $x=(1+\theta)/(2+\theta)$.
For $\omega<2+\theta$ equality also requires full order $R$ in every
contributing block; at the endpoint $\omega=2+\theta$ the scale does not
matter.
\end{proof}

For $\theta=1/2$,
$C_{1/2}=(25/6)\sqrt{5/3}$. Since
$M\geq(1+1/R)J_2/2$, the theorem gives
\begin{equation}\label{res:coordinate-bound}
 M\sqrt\nu\geq\frac{25}{12}\sqrt{\frac53}
       (1+1/R)\frac{n^{3/2}}{\sqrt R}.
\end{equation}
Divisible packings with $p+c=R$, $p:c=(1+\theta):1$ and
$H=bs/p$, $L=H/c$ attain~\eqref{res:holder-bound} with $bs$ in
place of $b(s-1)$. Their ratio to the lower bound is
$(s/(s-1))^{1+\theta}$. The cone-coordinate comparison has exactly the
same ratio because all blocks have full order. Thus these divisible
families attain the relaxation asymptotically as $s\to\infty$. When
$b\lceil s/(R-1)\rceil$ constrains the optimizer, or divisibility fails,
use~\eqref{res:joint-conditions} and the finite enumeration instead.

The same ratio $p:c$ follows directly by optimizing a uniform
full-order packing: $J_\omega H^\theta=(bs)^{1+\theta}R^\omega/(p^{1+\theta}c)$.
For static resources, $\theta=0$, the ratio is $1:1$; for the common
iteration-count factor $\sqrt\nu$ it is $3:2$. If $\omega>2+\theta$,
larger blocks increase $J_\omega H^\theta$, so the cap $R$ need not bind.
No ratio $p:c$ is universal once constraints on source width, minimum
block size, and integrality apply. In particular, $3:2$ need not be
optimal for a cost model based on dense cubic linear algebra, $\omega=3$.

Unshared PSD formulations have a different optimum. For divisible Schur groups of common width $g$ and $A=bs$,
\[
 L=\nu=A/g,\quad \rho=A(1+1/g),\quad
 M\sqrt\nu=A^{3/2}\frac{(g+1)(g+2)}{2g^{3/2}}.
\]
The first two resources times $\sqrt\nu$ decrease with $g$, whereas
the last has its unique integer minimum $15A^{3/2}/8$ at $g=4$, if this
width is available. Indeed, consecutive values decrease exactly when
$(g+1)^5-g^3(g+3)^2=-g^4+g^3+10g^2+5g+1\geq0$;
this holds for $g=1,2,3$ and fails thereafter. The polynomial is negative
at four and has negative derivative thereafter. Heterogeneous groups
cannot improve the value when $4\mid s$: at fixed mean width $u$, convexity
balances widths between two neighboring integers. The interpolated cost
is $au+b'$ with $a>0$ and $b'\leq0$, so the derivative of
$(au+b')u^{-3/2}$ changes sign, if at all, from positive to negative.
There is no interior minimum. Its minimum is therefore at an integer
width, where the preceding comparison applies. If four is unavailable,
the largest allowed width below four is optimal in the divisible model.
The exact integer counts above handle all remaining cases.
